\input{physics_report_base.tex}
\usepackage{multirow}
% ===== 数据输入区开始 =====
\SetField{experiment}{用扭摆法测定物体转动惯量}
% 后端以 \SetField{字段名}{值} 填入数据；缺省值保持空白。
% 母模板封面字段：\Field{experiment} \Field{student} \Field{student_id} \Field{group} \Field{date}
% ===== 数据输入区结束 =====
\begin{document}\fontsize{10.5pt}{20pt}\selectfont
\ReportCover{用扭摆法测定物体转动惯量}
\ReportSection{实验目的}
1. 观察小角度扭摆，测量空盘与规则物体的周期并标定扭转常数。\par
2. 比较圆柱、圆筒、球和细杆的理论与实验转动惯量。\par
3. 改变两个对称滑块到转轴的距离，验证平行轴定理及惯量对距离平方的依赖。\par
\ReportSection{实验原理}
扭摆的回复力矩在小角度范围内近似满足 $M=-\kappa\theta$，运动方程为 $I\ddot\theta+\kappa\theta=0$，故
\[T=2\pi\sqrt{\frac I\kappa}.\]
以圆柱作标准物体，比较空盘周期 $T_0$ 和装圆柱周期 $T_1$，可由圆柱理论惯量 $I'_1=\frac{mD^2}{8}$ 求 $\kappa=\frac{4\pi^2I'_1}{T_1^2-T_0^2}$。圆柱、圆筒置于载物盘上，实验惯量为 $I_j=\frac{\kappa T_j^2}{4\pi^2}-I_0$（$j=1,2$）；木球、细杆及细杆加滑块直接安装在转轴上，实验惯量为 $I_j=\frac{\kappa T_j^2}{4\pi^2}$，不扣除 $I_0$。对两个质量均为 $m_s$ 的滑块，平行轴模型为 $I(x)=I_{\rm rod}+2I_s+2m_sx^2$。
以 $\Delta=T_1^2-T_0^2$ 记周期平方差。暂把标准圆柱的理论惯量视为已知常量，独立重复计时的传播公式为
\[u_\kappa=\kappa\sqrt{\left(\frac{2T_0u_{T_0}}{\Delta}\right)^2+
\left(\frac{2T_1u_{T_1}}{\Delta}\right)^2}.\]
\[\begin{aligned}
u_{I_0}^2={}&\left(\frac{2I'_1T_0T_1^2u_{T_0}}{\Delta^2}\right)^2+
\left(\frac{2I'_1T_1T_0^2u_{T_1}}{\Delta^2}\right)^2,\\
u_{I_j}^2={}&\left(\frac{2I'_1T_ju_{T_j}}{\Delta}\right)^2+
\left(\frac{2I'_1T_0(T_j^2-T_1^2)u_{T_0}}{\Delta^2}\right)^2\\
&+\left(\frac{2I'_1T_1(T_j^2-T_0^2)u_{T_1}}{\Delta^2}\right)^2\quad(j=2).
\end{aligned}\]
木球、细杆及细杆加滑块的计时不确定度为
\[u_I^2=\left(\frac{\kappa T u_T}{2\pi^2}\right)^2+\left(\frac{T^2u_\kappa}{4\pi^2}\right)^2.\]
标定圆柱的 $I_1$ 与 $I'_1$ 由同一组周期构成，其计时项相消；完整不确定度还需计入圆柱质量、直径等仪器误差。
\begin{center}\includegraphics[width=3.4cm]{figures/torsion-apparatus.png}\par
{\small 扭摆装置示意图。}\end{center}
\ReportSection{实验内容}
\ReportSubsection{周期和几何量}
1. 调节底座水平与转轴竖直，保证载物盘空载时自由摆动；轻转一小角度后放手，避免向下压轴。\par
2. 测量空盘连续二十周期的总时长三次。圆柱、圆筒置于载物盘上；拆下载物盘后直接安装木球或细杆，各计十周期并重复三次。\par
3. 测量各物体质量及必要的直径、长度；装置应居中、固定可靠，换物体时保持转轴条件一致。\par
\ReportSubsection{平行轴验证}
4. 把两个滑块关于转轴对称放在细杆上，量取滑块质心到轴的距离；在 5、10、15、20 cm 四个位置分别计五周期三次。\par
5. 将实验惯量与理论惯量对照，并以 $x^2$ 为横轴作图，检查线性关系。
\ReportSection{实验仪器}
扭摆实验仪、电子天平、游标卡尺、刻度尺和秒表。按计时周期数换算单周期，空盘、规则物体与滑块的周期数不同，不可混用。
\clearpage
\ReportSection{7. 数据记录与处理}
1. 弹簧扭转常数和转动惯量的测定，数据记录见表 3.5.1。
\begin{center}
表 3.5.1\quad 刚体转动惯量测量实验数据记录表\par\smallskip
{\fontsize{9}{13}\selectfont\renewcommand{\arraystretch}{1.35}\setlength{\tabcolsep}{2pt}
\begin{tabular}{|m{1.35cm}|m{1.0cm}|m{1.05cm}|m{1.35cm}|m{.6cm}|m{1.35cm}|m{2.95cm}|m{2.85cm}|m{.85cm}|}\hline
\centering 物体名称 & \centering 质量 $m$\\（kg） & \multicolumn{2}{c|}{\shortstack{几何尺寸\\（$10^{-2}$ m）}} & \multicolumn{2}{c|}{周期（s）} & \centering 转动惯量\\理论值（kg$\cdot$m$^2$） & \centering 转动惯量\\实验值（kg$\cdot$m$^2$） & \centering 百分差\\$E_0$（\%）\tabularnewline\hline
\multirow{4}{1.35cm}{\centering 载物圆盘} & \multirow{4}{1.0cm}{\centering \begin{tikzpicture}[baseline=0pt]\draw (0,1.65) -- (.7,0);\end{tikzpicture}} & \multicolumn{2}{c|}{\multirow{4}{1.7cm}{\centering \begin{tikzpicture}[baseline=0pt]\draw (0,1.65) -- (1.6,0);\end{tikzpicture}}} & \multirow{3}{.6cm}{\centering $T_0$} & \Field[6mm]{period_disk_01} & \multirow{4}{2.95cm}{\centering \begin{tikzpicture}[baseline=0pt]\draw (0,1.65) -- (2.5,0);\end{tikzpicture}} & \multirow{2}{2.85cm}{\centering $I_0=I'_1\dfrac{\overline T_0^2}{\overline T_1^2-\overline T_0^2}$} & \multirow{4}{.85cm}{\centering \begin{tikzpicture}[baseline=0pt]\draw (0,1.65) -- (.75,0);\end{tikzpicture}}\tabularnewline\cline{6-6}
 &  & \multicolumn{2}{c|}{} &  & \Field[6mm]{period_disk_02} &  &  & \tabularnewline\cline{6-6}\cline{8-8}
 &  & \multicolumn{2}{c|}{} &  & \Field[6mm]{period_disk_03} &  & \multirow{2}{2.85cm}{\centering $\MathField{empty_inertia}$} & \tabularnewline\cline{5-6}
 &  & \multicolumn{2}{c|}{} & $\overline T_0$ & \Field[6mm]{period_disk} &  &  & \tabularnewline\hline
\multirow{4}{1.35cm}{\centering 实心圆柱体} & \multirow{4}{1.0cm}{\centering \Field[6mm]{mass_cylinder_kg}} & \multirow{3}{1.05cm}{\centering $D_1$} & \Field[6mm]{geometry_cylinder_01} & \multirow{3}{.6cm}{\centering $T_1$} & \Field[6mm]{period_cylinder_01} & \multirow{2}{2.95cm}{\centering $I'_1=\dfrac18m_1\overline D_1^2$} & \multirow{2}{2.85cm}{\centering $I_1=\dfrac K{4\pi^2}\overline T_1^2-I_0$} & \multirow{4}{.85cm}{\centering \Field[6mm]{error_cylinder}}\tabularnewline\cline{4-4}\cline{6-6}
 &  &  & \Field[6mm]{geometry_cylinder_02} &  & \Field[6mm]{period_cylinder_02} &  &  & \tabularnewline\cline{4-4}\cline{6-8}
 &  &  & \Field[6mm]{geometry_cylinder_03} &  & \Field[6mm]{period_cylinder_03} & \multirow{2}{2.95cm}{\centering $\MathField{theory_cylinder}$} & \multirow{2}{2.85cm}{\centering $\MathField{experiment_cylinder}$} & \tabularnewline\cline{3-6}
 &  & 平均值 & \Field[6mm]{diameter_cylinder_cm} & $\overline T_1$ & \Field[6mm]{period_cylinder} &  &  & \tabularnewline\hline
\multirow{8}{1.35cm}{\centering 空心金属圆筒} & \multirow{8}{1.0cm}{\centering \Field[6mm]{mass_tube_kg}} & \multirow{3}{1.05cm}{\centering $D_{\text{外}}$} & \Field[6mm]{geometry_tube_outer_01} & \multirow{3}{.6cm}{\centering $T_2$} & \Field[6mm]{period_tube_01} & \multirow{4}{2.95cm}{\centering $I'_2=\dfrac18m_2(\overline D_{\text{外}}^2+\overline D_{\text{内}}^2)$} & \multirow{4}{2.85cm}{\centering $I_2=\dfrac K{4\pi^2}\overline T_2^2-I_0$} & \multirow{8}{.85cm}{\centering \Field[6mm]{error_tube}}\tabularnewline\cline{4-4}\cline{6-6}
 &  &  & \Field[6mm]{geometry_tube_outer_02} &  & \Field[6mm]{period_tube_02} &  &  & \tabularnewline\cline{4-4}\cline{6-6}
 &  &  & \Field[6mm]{geometry_tube_outer_03} &  & \Field[6mm]{period_tube_03} &  &  & \tabularnewline\cline{3-4}\cline{5-6}
 &  & 平均值 & \Field[6mm]{diameter_tube_outer_cm} & \multirow{5}{.6cm}{\centering $\overline T_2$} & \multirow{5}{1.35cm}{\centering \Field[6mm]{period_tube}} &  &  & \tabularnewline\cline{3-4}\cline{7-8}
 &  & \multirow{3}{1.05cm}{\centering $D_{\text{内}}$} & \Field[6mm]{geometry_tube_inner_01} &  &  & \multirow{4}{2.95cm}{\centering $\MathField{theory_tube}$} & \multirow{4}{2.85cm}{\centering $\MathField{experiment_tube}$} & \tabularnewline\cline{4-4}
 &  &  & \Field[6mm]{geometry_tube_inner_02} &  &  &  &  & \tabularnewline\cline{4-4}
 &  &  & \Field[6mm]{geometry_tube_inner_03} &  &  &  &  & \tabularnewline\cline{3-4}
 &  & 平均值 & \Field[6mm]{diameter_tube_inner_cm} &  &  &  &  & \tabularnewline\hline
\multirow{4}{1.35cm}{\centering 木球} & \multirow{4}{1.0cm}{\centering \Field[6mm]{mass_ball_kg}} & \multirow{3}{1.05cm}{\centering $D_{\text{球}}$} & \Field[6mm]{geometry_ball_01} & \multirow{3}{.6cm}{\centering $T_3$} & \Field[6mm]{period_ball_01} & \multirow{2}{2.95cm}{\centering $I'_3=\dfrac1{10}m_3\overline D_{\text{球}}^2$} & \multirow{2}{2.85cm}{\centering $I_3=\dfrac K{4\pi^2}\overline T_3^2$} & \multirow{4}{.85cm}{\centering \Field[6mm]{error_ball}}\tabularnewline\cline{4-4}\cline{6-6}
 &  &  & \Field[6mm]{geometry_ball_02} &  & \Field[6mm]{period_ball_02} &  &  & \tabularnewline\cline{4-4}\cline{6-8}
 &  &  & \Field[6mm]{geometry_ball_03} &  & \Field[6mm]{period_ball_03} & \multirow{2}{2.95cm}{\centering $\MathField{theory_ball}$} & \multirow{2}{2.85cm}{\centering $\MathField{experiment_ball}$} & \tabularnewline\cline{3-6}
 &  & 平均值 & \Field[6mm]{diameter_ball_cm} & $\overline T_3$ & \Field[6mm]{period_ball} &  &  & \tabularnewline\hline
\multirow{4}{1.35cm}{\centering 金属细杆} & \multirow{4}{1.0cm}{\centering \Field[6mm]{mass_rod_kg}} & \multirow{3}{1.05cm}{\centering $l$} & \Field[6mm]{geometry_rod_01} & \multirow{3}{.6cm}{\centering $T_4$} & \Field[6mm]{period_rod_01} & \multirow{2}{2.95cm}{\centering $I'_4=\dfrac1{12}m_4\overline l^2$} & \multirow{2}{2.85cm}{\centering $I_4=\dfrac K{4\pi^2}\overline T_4^2$} & \multirow{4}{.85cm}{\centering \Field[6mm]{error_rod}}\tabularnewline\cline{4-4}\cline{6-6}
 &  &  & \Field[6mm]{geometry_rod_02} &  & \Field[6mm]{period_rod_02} &  &  & \tabularnewline\cline{4-4}\cline{6-8}
 &  &  & \Field[6mm]{geometry_rod_03} &  & \Field[6mm]{period_rod_03} & \multirow{2}{2.95cm}{\centering $\MathField{theory_rod}$} & \multirow{2}{2.85cm}{\centering $\MathField{experiment_rod}$} & \tabularnewline\cline{3-6}
 &  & 平均值 & \Field[6mm]{length_rod_cm} & $\overline T_4$ & \Field[6mm]{period_rod} &  &  & \tabularnewline\hline
\end{tabular}}
\end{center}
\noindent 计算弹簧扭转常数：
\[K=4\pi^2\frac{I'_1}{\overline T_1^2-\overline T_0^2}=\MathField{torsion_constant}\,\mathrm{N\,m}.\]
\clearpage
2. 转动惯量平行轴定理的验证，数据记录见表 3.5.2。滑块质量
$m_{\text{滑}}=\MathField{slider_mass_g}\,\mathrm g$（两个滑块的平均质量）。\par
表中 $I_5$ 为两个滑块绕通过滑块质心转轴的转动惯量
\[I_5=\frac18m_{\text{滑}}(D_{\text{内}1}^2+D_{\text{外}1}^2)+\frac16m_{\text{滑}}l_1^2.\]
式中，$m_{\text{滑}}$ 为滑块质量；$D_{\text{内}1}$、$D_{\text{外}1}$ 为滑块的内径和外径；$l_1$ 为其长度。\par
本次实验取\hfill $I_5=\MathField{slider_pair_inertia}\,\mathrm{kg\,m^2}$。\par
\begin{center}
表 3.5.2\quad 验证平行轴定理实验数据记录表\par\smallskip
{\small\renewcommand{\arraystretch}{1.6}\setlength{\tabcolsep}{5pt}
\begin{tabular}{|m{5.2cm}|*{4}{>{\centering\arraybackslash}m{1.65cm}|}}\hline
$x$（$10^{-2}$ m） & 5.00 & 10.00 & 15.00 & 20.00\tabularnewline\hline
\multirow{3}{5.2cm}{\centering $T$（s）} & \Field[6mm]{slider_period_5_01} & \Field[6mm]{slider_period_10_01} & \Field[6mm]{slider_period_15_01} & \Field[6mm]{slider_period_20_01}\tabularnewline\cline{2-5}
 & \Field[6mm]{slider_period_5_02} & \Field[6mm]{slider_period_10_02} & \Field[6mm]{slider_period_15_02} & \Field[6mm]{slider_period_20_02}\tabularnewline\cline{2-5}
 & \Field[6mm]{slider_period_5_03} & \Field[6mm]{slider_period_10_03} & \Field[6mm]{slider_period_15_03} & \Field[6mm]{slider_period_20_03}\tabularnewline\hline
$\overline T$（s） & $\MathField{slider_period_5}$ & $\MathField{slider_period_10}$ & $\MathField{slider_period_15}$ & $\MathField{slider_period_20}$\tabularnewline\hline
\shortstack[l]{实验值（kg$\cdot$m$^2$）\\[4pt]$I=\dfrac K{4\pi^2}\overline T^2$} & $\MathField{slider_experiment_5}$ & $\MathField{slider_experiment_10}$ & $\MathField{slider_experiment_15}$ & $\MathField{slider_experiment_20}$\tabularnewline\hline
\shortstack[l]{理论值（kg$\cdot$m$^2$）\\[4pt]$I'=I_4+2mx^2+I_5$} & $\MathField{slider_theory_5}$ & $\MathField{slider_theory_10}$ & $\MathField{slider_theory_15}$ & $\MathField{slider_theory_20}$\tabularnewline\hline
百分差 $E_0=\dfrac{|I'-I|}{I'}\times100\%$ & $\MathField{slider_error_5}$ & $\MathField{slider_error_10}$ & $\MathField{slider_error_15}$ & $\MathField{slider_error_20}$\tabularnewline\hline
\end{tabular}}
\end{center}
\ReportSection{平行轴定理的数据图}
横轴取滑块距离平方 $x^2$（cm$^2$），纵轴取总惯量的 $10^{-3}\,\mathrm{kg\,m^2}$ 数值。散点是周期换算的实验结果，直线连接对应理论值；两条序列均来自本次读数和输入的滑块质量。
\begin{center}\begin{tikzpicture}
\begin{axis}[width=.82\textwidth,height=4.8cm,xlabel={$\frac{x^2}{\mathrm{cm^2}}$},ylabel={$\frac{I}{10^{-3}\,\mathrm{kg\,m^2}}$},grid=both,legend pos=north west]
\addplot+[only marks,mark=*,mark size=2pt] coordinates {\PlotCoordinates{slider_graph_points}};
\addlegendentry{实测惯量}
\addplot+[no marks,thick,blue] coordinates {\PlotCoordinates{slider_graph_theory_points}};
\addlegendentry{平行轴理论}
\end{axis}\end{tikzpicture}\end{center}
\ReportSection{讨论与分析}
\ReportSubsection{图形与模型}
若平行轴定理适用，$I$ 对 $x^2$ 的点应近似落在直线上，斜率由两个滑块的质量决定。若图中随位置出现系统偏离，应检查距离是否从滑块质心量起、两个滑块是否对称，以及滑块本身惯量是否计入。\par
\ReportSubsection{重复计时的不确定度}
下列标准不确定度仅计重复计时贡献，未计质量、尺寸等仪器不确定度。
\begin{center}\small\begin{adjustbox}{max width=\textwidth}
\begin{tabular}{|c|c|c|c|c|c|}\hline
物体 & 载物圆盘 & 实心圆柱体 & 空心金属圆筒 & 木球 & 金属细杆\\\hline
$u_T$/s & $\MathField{period_disk_u}$ & $\MathField{period_cylinder_u}$ & $\MathField{period_tube_u}$ & $\MathField{period_ball_u}$ & $\MathField{period_rod_u}$\\\hline
$u_I$/kg m$^2$ & $\MathField{empty_inertia_u}$ & $\MathField{experiment_cylinder_u}$ & $\MathField{experiment_tube_u}$ & $\MathField{experiment_ball_u}$ & $\MathField{experiment_rod_u}$\\\hline
\end{tabular}\end{adjustbox}\end{center}
滑块四种位置的 $u_T$ 依次为
$\MathField{slider_period_5_u}$、$\MathField{slider_period_10_u}$、$\MathField{slider_period_15_u}$、$\MathField{slider_period_20_u}$ s；
$u_I$ 依次为 $\MathField{slider_experiment_5_u}$、$\MathField{slider_experiment_10_u}$、$\MathField{slider_experiment_15_u}$、$\MathField{slider_experiment_20_u}$ kg m$^2$。
\ReportSubsection{主要误差}
计时起止点不一致、转轴摩擦、振幅过大和物体偏心都会改变周期。圆柱用于标定 $\kappa$，其质量和直径误差会传递到所有实验惯量。仅圆柱、圆筒扣除空盘惯量 $I_0$；木球、细杆及细杆加滑块不扣除。\par
\ReportSection{实验结论}
由圆柱与空盘的周期差标定扭转常数 $\kappa=(\MathField{torsion_constant}\pm\MathField{torsion_constant_u})\,\mathrm{N\,m}$，再求其他物体惯量。这里的不确定度只计重复计时的贡献。滑块位置变化所得散点应与 $I'=I_4+I_5+2m_sx^2$ 的平方关系比较，不能只凭周期随距离变大作结论。
\ReportSection{思考题}
为何空盘计时二十周期，而物体只计十周期？空盘周期较短，扩大总计时长度可以减小启停反应时间的相对影响。为何使用两个对称滑块？对称布置可以减少质心偏离转轴引起的附加晃动。

\ReportRecord
% ===== 原始数据图片插入区 =====
\end{document}
